\begin{abstract}
University support offices can hold only a limited number of outreach meetings and must therefore
decide early which students to contact. This study examines whether a tabular foundation model,
\tabicl, should be used to rank students for such meetings after the first semester, and whether its
explanations are trustworthy and comprehensible. Using 3,630 Dropout or Graduate records from the
public UCI \emph{Predict Students' Dropout and Academic Success} dataset, all modelling choices were
fixed on a validation split before a single locked final test, in which each of five models flagged
the top 20\% of students, the assumed outreach capacity. \tabicl achieved a ROC-AUC of 0.957 and
flagged 109 test students, 108 of whom dropped out (precision 0.991, recall 0.507); however, it was
not reliably better than XGBoost (difference 0.001, 95\% interval $-0.006$ to 0.009) or an
interpretable Explainable Boosting Machine (0.003, $-0.007$ to 0.014). A post-hoc audit of all 544
validation students showed that the SHAP explanations were faithful to the model but often
unstable: when the comparison group changed, 411 students fell below the 0.80 stability threshold
for their top five reasons. Moreover, the flags largely restated information already visible to
advisers: 78 of the 108 flagged validation students had passed no first-semester units, and
first-semester progress was the leading reason for every flagged student. In a 50-person study,
participants predicted the model's flag equally well with and without explanation charts (0.900
versus 0.916; difference $-0.016$, 95\% interval $-0.060$ to 0.032). On this dataset,
these findings indicate that a simpler interpretable model would serve equally well, that
explanations should be presented with their uncertainty, and that the model can support, but not
replace, an adviser's judgement.
\end{abstract}

\noindent\textbf{Keywords:} student dropout risk, explainable artificial intelligence,
tabular foundation models, interpretable machine learning, SHAP, human-centred evaluation
